\scanpage{92}

En effet, $\overset{\inserted{30}}{27}$ est un cas particulier du th. 2, et ($\overset{\inserted{31}}{28}$) en dérive par passage à la limite pour $n\to+\infty$. \inserted{$*$} \deleted{En particulier, prenant} Pour un autre cas particulier du corollaire 1, \deleted{$f(s)=s^r$, $r>0$, on obtient} soient $p,q,\deleted{\unsure{r}}s$ tels que

% margin note written sideways

\inserted{Prenons en particulier $f(\unsure{u|s})\equiv1$, on obtient alors $\boxed{\text{(32)}\quad\sum\rho_i(uv)\leq\sum\rho_i(u)\rho_i(v)}$. Tenant compte du \deleted{(21)} \inserted{N\textsuperscript{o} 1, th. 3}, \unsure{pour} $p=1$, \deleted{et \unsure{soit}} on en conclut que si le deuxième mb est fini, alors $uv$ est un op. de Fredholm. \deleted{Par \unsure{conséquent}} Or si \deleted{\unsure{$\sum\lambda_i$}} \inserted{$E=G$}, on a \deleted{d'où $|\mathrm{Tr}(uv)|\leq\sum|\lambda_i(uv)|\leq\sum\rho_i(uv)$} $|\mathrm{Tr}\,uv|\leq\deleted{\unsure{\sum\rho_i(v)|\sum}}\inserted{\sum\rho_i(uv)}$ \deleted{(inégalité \unsure{fond.} (21), pour $p=1$) d'où $|\mathrm{Tr}(uv)|\leq\sum\rho_i(u)\rho_i(v)$. Par suite on a} (N\textsuperscript{o} 1, th. \unsure{3|2}). \uline{Corollaire \unsure{2|3}}. Soient $E,F$ deux espaces de Hilbert, $u$ une application lin. compacte de $E$ dans $F$, $v$ une application lin. compacte de $F$ dans \unsure{$G$|$E$}, supposons $\sum\rho_i(u)\rho_i(v)<+\infty$, alors $vu$ est un op. de Fredholm, et on a $\boxed{\text{(33)}\quad|\mathrm{Tr}\,uv|\leq\sum\rho_i(u)\rho_i(v)}$ \deleted{\illegible{3}}. Cet énoncé précise que le produit de deux op. de H.-S. est un op. \deleted{\unsure{est}} de Fredholm. Généraliser l'inégal. du (29).}

\[
\frac1p+\frac1q=\frac1s\qquad(p,q,s\geq0)\tag*{(34)}
\]

posons $f(s)=s^{\unsure{s|r}}$, \deleted{on a donc \unsure{en le corollaire 1 donne|dans le corollaire 1}} l'inégalité \deleted{du \unsure{cor.}} (\unsure{31|34}) donne

\[
\left(\sum_{i=1}^{n}\rho_i(uv)^s\right)^{1/s}\leq\left(\sum_{i=1}^{n}(\rho_i(u)\rho_i(v))^s\right)^{1/s}
\]

Or l'inégalité de Hölder \deleted{\unsure{$\sum,\sum$} s'applique \unsure{ici que le deuxième membre puis} $\sum(\rho_i(u))^s(\rho_i(v))^s\leq$} le deuxième mb est majoré par $\left(\sum\rho_i(u)^p\right)^{1/p}\left(\sum\rho_i(v)^q\right)^{1/q}$.\ \unsure{Par suite|On trouve}:

\uline{Corollaire 2}. (``\uline{Inégalité de Hölder \deleted{pour} opérationnelle}''). \deleted{Si $u$ \unsure{si $u$ est}} Soient $E,F,G$ trois espaces de Hilbert, soit $u$ un op. cpt de $E$ dans $F$, $v$ un op. cpt de $F$ dans $G$, alors on a, pour tout \deleted{\unsure{entier} $\geq1$} $(p,q,s)$ \deleted{\unsure{et tout} $p$} satisfaisant à (34), et tout entier \unsure{$n>1$|$n\geq1$|$n>0$},

\[
\left(\sum_{i=1}^{n}\rho_i(uv)^{\unsure{r|s}}\right)^{\unsure{1/r|1/s}}\leq
\left(\sum_1^n\rho_i(u)^p\right)^{1/p}\left(\sum_1^n\rho_i(v)^q\right)^{1/q}\tag*{(35)}
\]

\deleted{En particulier}

\[
\deleted{\left(\sum_i\rho_i(uv)^s\right)^{1/s}\leq
\left(\sum_i\rho_i(u)^p\right)^{1/p}\left(\sum_i\rho_i(v)^q\right)^{1/q}}
\tag*{\deleted{(36)}}
\]

\[
\deleted{\unsure{\mathcal N_r(uv)|N_r(uv)|N_s(vu)}\leq \mathcal N_p(u)\mathcal N_q(v)}
\quad \inserted{\|uv\|_r\leq\|u\|_p\|v\|_q}\quad
\left(\frac1r=\frac1p+\frac1q\right)
\]

\deleted{(et on a \unsure{tout aussi lieu} si tous les \unsure{opérateurs})}

Le cas le plus important de l'inégalité (37) est celui où $s=1$, i.e. où $p$ et $q$ sont deux ``exposants conjugués''; \deleted{Il est \unsure{que mieux dans ce cas|nous y reviendrons}} \inserted{dans ce cas, (36) est d'ailleurs contenu dans} l'\deleted{inégalité du N\textsuperscript{o} suiv., en une formule plus frappante.} \unsure{(38)}

\[
\boxed{\|uv\|_1\leq\|u\|_p\|v\|_{p'}\quad \left(\frac1p+\frac1{p'}=1\right)}\tag*{(38)}
\]

\deleted{(39) \unsure{si tous} : l'inégalité de Hölder ci-dess\unsure{us}}. \inserted{(38) résulte d'ailleurs aussi de (32), grâce à l'inégalité \unsure{de Hölder ordinaire}}

% slanted

\margin{\unsure{Expliquer|Expliq.} \unsure{l'obtention|l'exécution} \unsure{d'ordre|d'un} \illegible{1}}

\uline{Corollaire 4}. Soient $E,F,G$ trois espaces de Hilbert de \unsure{même} dimension finie $n$, $u$ un op. cpt de $E$ dans $F$, $v$ un op. compact de $F$ dans $G$, $f$ étant une fonction \deleted{\unsure{sym.}} de $n$ arguments $t_i>0$ telle que $f(e^{t_i},\dots e^{t_n})$ soit fonction convexe de $(t_i)\in R^n$. Sous ces conditions, si \deleted{\unsure{l'un de ces}} $u$ et $v$ sont inversibles on a \deleted{la \unsure{formule}} \unsure{si $f$ est} \unsure{définie et continue} pour $s_i\geq0$, on a \unsure{la limite}

\margin{?}

\[
f(\rho_1(uv),\ldots,\rho_n(uv))\leq f(\rho_1(u)\rho_1(v),\ldots,\rho_n(u)\unsure{\rho_n(v)})
\tag*{(39)}
\]

En effet, \deleted{\unsure{on sait du}} \inserted{utilisant les conditions} du th. 2\textsuperscript{o}, hyp. b), on

% slanted

\margin{\uline{petits} \uline{caractères}}

a en vertu de

\[
\prod\rho_i(uv)=\deleted{\unsure{f(uv)|f(vu)}}\prod\rho_i(u)\rho_i(v)\tag*{(40)}
\]

\deleted{En effet,} Pour vérifier (40), on \deleted{doit \unsure{envisager si on a|élève au carré}} \deleted{est ramené} \inserted{peut supposer} que $E=F=G$, alors les deux mb. de (35) sont resp\textsuperscript{t}. $|\det(uv)|$ et $|\det u|\,|\det v|$ (formule (16)) d'où le résultat. Si $u$ et $v$ sont q.qcon\unsure{ques} mais $f$ définie et continue pour $t_i\geq0$, (39) sera encore valable par prolongement par continuité à partir du cas précédent.

Donner des \unsure{liaisons|variations|variantes} sur la composition des opérateurs comme multiplication

% archive sheet number: 93
